% =============================================================================
% Chapter 6 --- Results
% =============================================================================

\chapter{Results}
\label{ch:results}

This chapter consolidates the simulation evidence and interprets it
against the architectural intent of the design. The presentation is
organized around the three claims of \Cref{sec:verif-goals} --- the
static transaction-word contract, the dynamic FSM protocol, and the
closed-loop recovery policy. The directed regression set is summarized
quantitatively in \Cref{tab:results-summary}, and the two most
architecturally informative experiments --- the full closed-loop
program--verify--erase recovery and the inference dataflow --- are
illustrated in
\Cref{fig:results-recovery,fig:results-inference}. Per-test waveform
captures and console logs for every entry in \Cref{tab:results-summary}
are archived in the project repository.

\section{Regression Coverage and Outcomes}
\label{sec:results-summary}

\Cref{tab:results-summary} summarizes every directed test in the
regression set, grouped by the architectural claim it discharges and
annotated with the number of stimulus cases it covers and the resulting
pass/fail outcome. All 60+ directed cases pass with zero scoreboard
mismatches.

\begin{table}[h]
    \centering
    \caption{Directed regression coverage. Each row corresponds to a
        focused test of one architectural claim. \emph{Cases} counts the
        number of distinct stimuli inside the bench (e.g.\ randomized
        address indices, recovery scenarios, pixel windows). The two
        ``hero'' regressions are highlighted; the rest are summarized in
        a single row per testbench.}
    \label{tab:results-summary}
    \begin{tabularx}{\linewidth}{l X c c}
        \toprule
        \textbf{Test} & \textbf{Architectural claim exercised} &
        \textbf{Cases} & \textbf{Outcome} \\
        \midrule
        Address \& pulse-mode packing &
            Static contract: outbound core word matches the
            host-requested cell, and the pulse-mode bus carries the
            encoding appropriate to the command class. Randomized PE,
            SA, column, row, and payload byte. &
            30 & PASS \\
        PI / Controller / Buffer integration &
            Dynamic protocol of the three-module stack under mixed
            command traffic, with idle gating between transactions
            (host reads, programs, erases, and a full inference burst). &
            10 & PASS \\
        Host read reorder &
            Stability of the captured-address register and
            address-to-core-word reconstruction under out-of-program
            access order. Eight cells, fixed reordered read schedule. &
            8 & PASS \\
        Host erase selectivity &
            Spatial selectivity of host-initiated erase: erased cell is
            cleared while neighbor at a distinct address is preserved. &
            2 & PASS \\
        \textbf{Inference dataflow (hero)} &
            Collect $\rightarrow$ compute $\rightarrow$ result flow,
            bit-serial output lanes, and bit-select advance from
            zero to seven. Sanity check: compute must not start before
            eight pixel writes. &
            1 (9 beats) & PASS \\
        10-weight program--verify &
            Compact regression with all four per-cell recovery scenarios:
            match, under (retry), over (erase), retry-exhausted (fall to
            erase). Suitable for GUI waveform debug. &
            10 & PASS \\
        \textbf{Full program--verify regression (hero)} &
            Closed-loop recovery policy at scale, with deliberate
            under- and over-programmed fault injection at disjoint
            pre-declared index sets spanning the whole run. End-of-run
            shadow-matrix scoreboard $+$ branch-exercise discipline. &
            $\ge$ 640 cells & PASS \\
        Unit benches (buffer $\times$ 3, parallel interface) &
            Per-module contracts in isolation: one-hot validity,
            reset-to-zero buffer, full-capacity overwrite, bit-serial
            output. &
            4 & PASS \\
        \bottomrule
    \end{tabularx}
\end{table}

The remainder of the chapter interprets these outcomes against the
architectural claims.

\section{Static Contract: Address and Pulse-Mode Packing}
\label{sec:results-static}

The 30-case address-and-pulse-mode regression sweeps randomized cell
coordinates and payload bytes through program, read, and erase commands.
For each case the testbench reconstructs the expected outbound core word
and pulse-mode encoding from the same helper functions used by the RTL,
and compares them against the live outputs bit-for-bit throughout the
active phase. Every case passes. This validates two things at once: the
bidirectional addressing round-trip (host one-hot tail $\rightarrow$
packed binary $\rightarrow$ rebuilt one-hot tail) is implemented
consistently across the interface and the controller, and the
command-to-pulse-mode mapping is correctly wired through the pulse-output
combinational block. Static-contract regressions are the floor on which
every dynamic claim subsequently rests; their cleanliness is what makes
FSM-level conclusions interpretable.

\section{Dynamic Protocol: Three-Module Integration}
\label{sec:results-integration}

The integration regression exercises the three-module stack with ten
labelled scenarios that span host reads, host programs, host erases, and
inference, with idle gating between transactions. Each program command
leaves the addressed buffer slot updated with the host-supplied payload,
reflecting the data-load write that the program sub-FSM performs in its
\texttt{SELECT} state. Read and erase commands leave the buffer
untouched. The nine-beat inference burst loads exactly eight pixels into
the buffer before transitioning to compute, and the bit-serial lanes
expose the expected per-bit pattern during the compute window.

What this regression establishes is the consistency of the dynamic
protocol \emph{across} module boundaries: there is no missed handshake
between the interface and the controller, no contention between the
controller's buffer-mode selection and the buffer's mode-qualified ports,
and no state-machine residue from one transaction leaking into the next.
The host read reorder, host erase selectivity, and unit benches further
narrow this claim down to specific design risks (stale captured-address
register, mis-routed erase decode, per-module contracts) and all return
clean. With the static and dynamic regressions both clean, the
closed-loop recovery claim is the only remaining architectural claim
that needs separate evidence.

\section{Closed-Loop Recovery (Hero Result)}
\label{sec:results-recovery}

The full program--verify regression is the central architectural result
of the project. It programs an entire weight matrix sourced from a memory
image; during the run, a set of pre-declared cell indices is forced
under-programmed in the verify window and a disjoint set is forced
over-programmed (\Cref{sec:verif-injection}). Both sets span the entire
run, exercising the recovery policy from start to finish. End-of-run
checks combine a shadow-matrix scoreboard (every cell must converge to
its expected weight) with a branch-exercise discipline (both recovery
counters must be non-zero).

\paragraph{Outcome.}
Every cell converges. The shadow-matrix scoreboard reports zero
mismatches. The under-programmed recovery counter shows a substantial
number of re-program events, including indices that hit the retry budget
and contribute erase events; the over-programmed counter shows the
expected smaller number of erase-driven recoveries. Neither is zero, so
the fault-injection gating is honored.

A regression that programs an entire weight matrix in the presence of
deliberately injected verify failures, and still converges every cell to
the expected value, cannot be passing by accident: the recovery policy
must be entering the right branch on every failure, the retry pulse
train must be applying meaningful corrections on every retry, and the
address bookkeeping must be returning to the correct cell on every
re-entry into the program state. The clean scoreboard is the strongest
piece of evidence the project produces, and it is the result that
justifies handing the controller to the analog phase.

\begin{figure}[h]
    \centering
    \includegraphics[width=\textwidth]{example-image.png}
    \caption{Closed-loop recovery, hero waveform from the full
        program--verify regression. The capture spans one
        over-programmed injection followed by one under-programmed
        injection, end to end. \emph{Top:} an over-programmed
        \texttt{VERIFY\_CHECK} routes the FSM into \texttt{ERASE} with
        the verify-initiated flag; the erase sub-state advances through
        its five phases; \texttt{ERASE\_COMPLETE} returns to
        \texttt{PROGRAM}; the cell is reprogrammed and successfully
        verified on the next attempt. \emph{Bottom:} an under-programmed
        \texttt{VERIFY\_CHECK} increments the retry counter and re-enters
        \texttt{PROGRAM} with the short single-cycle nudge train; the
        next verify reads a match and the FSM returns to idle. Together
        the two segments exercise the full
        \texttt{PROG}~$\rightarrow$~\texttt{VERIFY}~$\rightarrow$
        \texttt{ERASE}~$\rightarrow$~\texttt{RE-PROG}~$\rightarrow$
        \texttt{VERIFY} closed-loop path defined in
        \Cref{sec:rtl-controller-recovery}.}
    \label{fig:results-recovery}
\end{figure}

\paragraph{Asymmetric injection.}
The under-programmed set is deliberately larger than the over-programmed
set, reflecting the device-physics observation that under-programming is
the more common analog failure mode --- a single program pulse undershoots
more often than it overshoots. The regression therefore spends more
time exercising the retry path than the erase path, in proportion to the
expected real-world distribution.

\section{Inference Dataflow (Hero Result)}
\label{sec:results-inference}

The inference regression validates the bit-serial output discipline of
the input buffer and the cycle-by-cycle advance of the bit-select line
in the controller --- the two pieces of the design that the analog MVM
array depends on for its per-bit-time input vector. A single misalignment
between the bit-select advance and the buffer's lane outputs would
corrupt the entire MVM result, so this regression is the structural
analogue of the recovery hero for the read-only path.

A nine-beat inference burst drives the controller through collect,
compute, and result. Eight pixel writes are observed in the buffer
before any compute indicator fires, the bit-select counter steps from
zero to seven during compute, and the eight bit-serial lanes present the
selected bit of each of the eight buffered pixels in parallel.
\Cref{fig:results-inference} captures the alignment.

\begin{figure}[h]
    \centering
    \includegraphics[width=\textwidth]{example-image.png}
    \caption{Inference dataflow, hero waveform. \emph{Collect phase:}
        eight successive data-load cycles write pixels into the first
        row of the buffer register array. \emph{Compute phase:} the
        pulse-mode bus carries the inference encoding, the bit-select
        line steps from zero to seven across consecutive cycles, and the
        eight bit-serial lanes present the corresponding bit of each
        pixel in parallel. \emph{Result phase:} a one-cycle wrap-up
        fires on pulse completion before the FSM returns to idle. The
        alignment between \emph{bit\_sel} and the lane outputs is the
        invariant on which the MVM result depends.}
    \label{fig:results-inference}
\end{figure}

\section{What the Results Establish and What They Do Not}
\label{sec:results-scope}

The collected evidence establishes, under explicit modeling assumptions,
that the controller's digital behavior conforms to its architectural
intent across every claim of \Cref{sec:verif-goals}. The static contract
is honored on every directed case; the dynamic protocol is honored on
every integration scenario; the closed-loop recovery policy converges
every weight in the full regression to its expected value in the
presence of deliberately injected failures; and the inference dataflow
correctly stages pixels and presents the expected per-bit values during
compute.

Equally important to be explicit about what is \emph{not} established.
Timing closure, physical feasibility, and gate-level power/area are out
of scope --- no synthesis or place-and-route is in the verification
flow. End-to-end numerical correctness of an analog MVM is not
established either, because no golden MAC reference exists in the
environment. And the controller's behavior under analog non-idealities
not captured by the four-bit quantized shadow matrix --- in particular,
the timing of the \emph{op\_done} handshake relative to actual cell
settling, and the cell-level variability of the Y-Flash device under
repeated programming --- remains a question for the analog phase, not
for the digital phase reported here. These scope limits motivate the
follow-on roadmap of \Cref{ch:limitations}.
